\section*{Resolução da Questão 6: Busca Guiada por Monte Carlo (MCTS)}

\subsection*{A) Mecanismo e Pilares do Monte Carlo Tree Search}

\begin{enumerate}[label=\alph*)]
    \item \textbf{Seleção Estratégica:} A CPU escorrega pelas ramificações existentes avaliando matematicamente (UCB1) os pesos de risco/recompensa, até esbarrar no fim de um ramo inexplorado.
    \item \textbf{Bifurcação (Expansion):} Acopla-se uma nova possibilidade inédita de jogada nesse limite da árvore, destravando um braço cego para ser vasculhado.
    \item \textbf{Simulação Desenfreada (Rollout):} Daquele ponto em diante, lança-se uma chuva de lances fictícios quase cego (ou moderado por heurística), testando exaustivamente o tabuleiro até presenciar uma vitória de algum dos lados ou empate.
    \item \textbf{Eco do Sucesso (Backpropagation):} A mensagem do final do jogo "eco" flui ladeira acima pela árvore inteira, incrementando os termômetros de acerto ($W$) e uso ($N$) em cada pai no percurso.
    \item \textbf{Influência da Confiança ($N$):} O eixo divisor da fórmula. Nós hiper-testados exalam segurança empírica, derrubando severamente a sua pontuação inflacionada na equação de UCT para que a CPU pare de testá-los e vá olhar outros cantos obscuros.
    \item \textbf{Métricas de Retorno ($W$):} A base tangível do acerto. É o medidor pragmático de qualidade do braço (o foco puro de "explotação").
    \item \textbf{O Dilema Explotação vs Exploração:} Um braço de ferro em que a Exploração sacrifica lucros imediatos para iluminar cantos esquecidos da árvore (na caça por rotas geniais escondidas), contra a Explotação que suga avidamente o que já tem garantia provada estatística de dar certo. A constante $C$ atua como o juiz regulador das duas.
\end{enumerate}

\subsection*{B) Execução Rastreável (10 Pulsos)}

Com calibragem $C=1.4$ e lances randomizados:

\begin{table}[H]
\centering
\begin{tabular}{|c|c|c|c|c|}
\hline
\textbf{Pulso} & \textbf{Ramificação Adotada} & \textbf{Vértice Criado} & \textbf{Resultado} & \textbf{Ajuste dos Pesos} \\ \hline
1 & c1 & c1 & Empatou & N=1, W=0.5 \\ \hline
2 & c2 & c2 & Amarelo ganhou & N=1, W=0.0 \\ \hline
3 & c3 & c3 & Empatou & N=1, W=0.5 \\ \hline
4 & c4 & c4 & Vermelho massacrou & N=1, W=1.0 \\ \hline
5 & c4 $\rightarrow$ c1 & c1 & Vermelho massacrou & N=1, W=0.0* \\ \hline
6 & c1 $\rightarrow$ c1 & c1 & Empatou & N=1, W=0.5 \\ \hline
7 & c3 $\rightarrow$ c1 & c1 & Amarelo ganhou & N=1, W=1.0 \\ \hline
8 & c4 $\rightarrow$ c2 & c2 & Vermelho massacrou & N=1, W=0.0* \\ \hline
9 & c4 $\rightarrow$ c3 & c3 & Vermelho massacrou & N=1, W=0.0* \\ \hline
10 & c2 $\rightarrow$ c1 & c1 & Empatou & N=1, W=0.5 \\ \hline
\end{tabular}
\end{table}

\subsection*{C) Cômputo Rígido do UCT}

A calibragem exige a balança de exploração:
\[ Formula \,\, UCT = \left(\frac{w_i}{n_i}\right) + C \times \sqrt{\frac{\ln N}{n_i}} \]

Para um painel genérico de $N = 3 + 5 + 1 = 9$ trânsitos da raiz:

\begin{itemize}
    \item \textbf{Eixo 1 (c1):} $0.666 + (1.4 \times 0.855) = 1.864$
    \item \textbf{Eixo 2 (c2):} $0.800 + (1.4 \times 0.662) = 1.727$
    \item \textbf{Eixo 5 (c5):} $1.000 + (1.4 \times 1.482) = \textbf{3.075}$
\end{itemize}

O eixo **c5** absorveu a liderança. Apesar de c2 ter alta fatia real de vitórias brutas (explotação de $0.80$), ele foi visitado um excesso de vezes (punido pela raiz de log sobre $n_i$). c5 possui incerteza máxima, e a variável de bônus exploratória injetou gigantescos $2.075$ em sua veia, arrastando-o para as prioridades.

\subsection*{D) Experimento: Distorção da Personalidade da IA (C)}

Sintetizamos 500 ciclos para observar a polarização da máquina.

\begin{itemize}
    \item \textbf{Personalidade Focada ($C=0.2$):} Pulverização de nós foi $c1=21, c2=6, c3=19, c4=454$. Um vício absurdamente agressivo num só pilar que se provou levemente positivo nos primeiros instantes. Se fosse falso positivo, a máquina travaria rumo a um colapso medíocre sem margem para recuperar o prejuízo.
    \item \textbf{Personalidade Curiosa ($C=5.0$):} A pulverização foi radicalmente insana: $c1=111, c2=125, c3=130, c4=134$. A máquina não consegue confiar nos próprios resultados estelares, duvidando de seu aprendizado e revisitando todos os buracos continuamente, perdendo vasto poder de processamento em eixos podres por pânico exploratório.
\end{itemize}

\subsection*{E) O Peso do Rollout Heurístico}

A Injeção de "malícia lógica" no motor de Rollout:

\begin{table}[H]
\centering
\begin{tabular}{|p{3cm}|p{6cm}|p{6cm}|}
\hline
\textbf{Sintoma} & \textbf{Aleatoriedade Cega} & \textbf{Semi-Guloso Guiado} \\ \hline
\textbf{Aceleração de HW} & Suprema. Sem gargalos de "IFs" verificadores de tabuleiro, gera trilhões de nós/segundo. & Freada substancial de hardware por injetar if/else de heurísticas centrais em cada turno simulado. \\ \hline
\textbf{Inteligência Estatística} & Letárgica. Aprende a não cometer suicídio perdendo centenas de peças primeiro até que os blocos de W absorvam o desastre. & Veloz e Afiada. Cada amostra de simulação provê dados puros e altamente factíveis com lances razoáveis de combate humano, não-poluídos. \\ \hline
\end{tabular}
\end{table}
